\section{Justice for the Wronged, Mercy for the Penitent}

Those who have suffered wrong can await vindication from Christ. Augustine describes the innocent person condemned and the guilty person released (\emph{City of God} XX.2). Human evidence can be hidden, institutions can fail, and a victim can die without public vindication. At the last judgment God's justice will be manifest, including where an earthly court failed. The truth of the injury remains known to him even when it can no longer be established before human witnesses.

This promise honors the truth of suffering. A Christian should never be required to call a lie true, an abuse harmless, or a betrayal good in order to express trust in God. Providence can bring good from what creatures do wickedly; the creature's wickedness remains wickedness. The Judge knows what was taken and who was injured. He knows the effects which could not be put into words and the losses which no earthly reparation could wholly restore. Hope in his judgment allows the wounded person to entrust the unfinished claim to someone who understands it completely.

That entrusting does not remove the duties of the present. The vulnerable must be protected, truthful testimony given where required, and just remedies pursued. Forgiveness does not require leaving a person exposed to further injury or pretending that trust has already been restored. Nor does recourse to legitimate justice itself mean a refusal to forgive. We can renounce personal vengeance while seeking protection and accountability. The eternal judgment belongs to Christ; the responsibilities already placed in our hands remain ours to fulfill.

The one who has done wrong must hear the same doctrine from the other side. A hidden injury does not cease to matter because its victim is silent. A person who has been forgiven cannot treat forgiveness as a reason to refuse restitution. Christ's mercy reaches the actual sinner and therefore reaches the actual wrong. It gives courage to name what happened without changing its description to preserve self-esteem. A truthful confession can be the beginning of a more demanding love than the evasive peace we had maintained before it.

St.~Thomas locates mercy at the root of divine justice. Whatever is due to a creature rests upon a prior gift: the nature, existence, and capacities which God freely bestowed. Tracing these gifts to their source reaches the goodness of the divine will (ST I, q.~21, a.~4). Even a creature's claim to what is due presupposes God's generosity. Justice and mercy are therefore present together in his works. When we ask him for mercy, we appeal to the goodness already expressed in our creation and every gift by which he has sustained us.

Jesus Christ gives himself for sinners while revealing the seriousness of the sin from which he saves. Forgiveness concerns a real wrong. It opens the sinner's return, heals the will, and calls the penitent to repair his relation with the person injured. This can be costly: a confession may require relinquishing an advantage maintained by falsehood. The victim can seek justice, and the wrongdoer can seek forgiveness, from the Savior who knows what happened to each.

The final disclosure also concerns the grace which delivered sinners. In the account transmitted in the posthumously compiled Supplement to the \emph{Summa}, Thomas holds that forgiven sins will be made known, together with the repentance and mercy by which the saints were saved (Supplement, q.~87, a.~2). He calls this the more probable and more common opinion and records Peter Lombard's contrary view that sins erased by repentance will not be disclosed to others. Thomas adds that such disclosure would bring the blessed no shame. God's justice and providence will be manifest; Catholic theologians differ over this particular manner of manifestation.

In either account, a sin truly remitted remains forgiven. The Church already remembers Peter's denial together with his repentance and faithful witness. Knowledge of the fall helps us understand what he received from Christ and why he owed him gratitude. Thomas's explanation of final disclosure concerns this acknowledgment of saving mercy, with the absence of shame which he expressly affirms. It gives no reason to imagine the blessed exposed to the cruelty with which people sometimes receive another person's faults now.

Meditation on that day can therefore purify our use of information now. We may need to know an accusation to protect someone or discharge a duty. We do not need to enjoy another person's disgrace. We may need to correct a false account without making a life reducible to its worst act. The Lord's judgment frees us to be exact about wrong while humble about the person. We too depend upon mercy, and we too have a history whose full meaning exceeds what our neighbors can see.

An enemy's conversion is therefore something we can ask for, including when the injury still needs a just remedy. If the wrong was ours, we can seek forgiveness and undertake the repair owed. A truthful account of the act is necessary to both responses. Prayer before judgment can help us acknowledge that act without making a claim about the person's eternal destiny, whether the person is an enemy or ourselves.

